\subsection{Surrogate model from the flow simulations and the question of grading}\label{sec:exchangersurrogate}

\paragraph{Why a surrogate?} For every design, the preheater calculation needs Nu and $f$ on both sides. So far there
were three ways to get them, each with a catch:
\begin{enumerate}
\item \textbf{Correlations from the literature} are fast but scatter widely. At Re 3000, Renon, Li 2020 and Li 2022
  give $f$ = 0.54, 0.15 and 0.81 for the same gyroid, more than a factor of 5, and Nu = 75, 63 and 79
  \cite{renon2024,alsafadi2026}. They also apply only to the cells and walls they were measured or computed for.
\item \textbf{Our own flow simulations} match our geometry exactly but cost 1--70 minutes per point. The segment model
  below evaluates about 440,000 combinations of approach, velocity, passes, cell and segment per variant. At two
  minutes each, that would take almost two years.
\item \textbf{Fitting our own correlation} (power law $\mathrm{Nu} = c\,\mathrm{Re}^m\Pr^n$) to the computed points
  would be fast, but it fixes the form in advance and does not say how reliable it is between the points.
\end{enumerate}
A \emph{surrogate} combines the advantages. It is trained once on a few expensive runs and afterwards delivers values
for any input in microseconds, together with a measure of how reliable they are. The concrete question here is whether
grading the cells along the flow path pays off. The literature only gives a blanket factor for that (Nu $\times$ 1.28
\cite{oh2025}). To check it ourselves, we need to know how Nu and $f$ depend on cell size, and that is exactly what the
surrogate learns from our own runs.

\paragraph{How a Gaussian process works.} The surrogate is a Gaussian process (kriging). It assumes no formula, only that
similar inputs give similar results. A kernel describes this similarity:
\[
  k(x,x') = \sigma^2 \exp\Big(-\sum_i \frac{(x_i-x_i')^2}{2\ell_i^2}\Big) + \sigma_n^2\,\delta_{xx'}.
\]
The length scale $\ell_i$ says over what distance in input $i$ the result changes noticeably. A large $\ell_i$ means
input $i$ hardly matters. $\sigma_n^2$ absorbs the noise of the runs, such as the time fluctuation of the mixed-mean
temperature. All parameters are determined from the data (maximum marginal likelihood). At a new point the process
returns a mean, a weighted blend of the known points, and a standard deviation. The latter is small near computed points
and grows in gaps, so one can see where the model is guessing. The length scales are a result in their own right: they
show which input really matters.

\paragraph{Why multi-level?} Fine grids are accurate but expensive (n144: 30--70 minutes, 8\,GB of graphics memory);
coarse grids are cheap (n48: 1--2 minutes) but systematically too high. A multi-level surrogate (co-kriging
\cite{kennedy2000}) uses both. Many coarse runs describe the \emph{shape} of the dependence on Re, Pr and cell. A few
fine runs determine the \emph{shift} caused by the grid. This works well if the grid error varies slowly and smoothly
across the input space, and that is checked below.

\paragraph{Design of experiments.} For each of gyroid and diamond, 14 points were chosen by Latin hypercube
\cite{mckay1979}: Re 500--10,000 (logarithmic), Pr 1.2--2.0, cell size $a$ = 0.75--1.6 times the standard cell, wall
fixed at 0.5\,mm. These points ran on the coarse grid (48 cells per $d_h$). Three control points ran on the medium grid
(96), and one run per type on the fine grid (144 cells per $d_h$, 62.5 million cells). Together with the runs of
Section~\ref{sec:exchangercfd} that makes 40 runs. The series ran unattended from a job queue that could be paused at any
time. A second AI agent watched memory, disk and divergence. No run failed.

\paragraph{Structure.} The data are referred to the correlation of Renon \cite{renon2024}, so that the Gaussian process
only has to learn the deviation:
\[
  y_\mathrm{Nu} = \ln\frac{\mathrm{Nu}}{\mathrm{Nu}_\mathrm{Renon}(\mathrm{Re},\Pr)},\qquad
  y_f = \ln\frac{f}{f_\mathrm{Renon}(\mathrm{Re})},\qquad x = (\ln\mathrm{Re},\ \Pr,\ a).
\]
\begin{enumerate}
\item \textbf{Coarse level:} per type, a Gaussian process with an anisotropic RBF kernel and a noise term over the 15
  points on 48 cells. Leave-one-out gives a mean error of 2.8\,\% (gyroid) and 0.9\,\% (diamond) for Nu. For $f$ it is
  9.5\,\% and 1.1\,\%: the gyroid friction scatters noticeably.
\item \textbf{Grid correction:} at each control point, $d_{96} = y_{96} - y_{48}(x)$. Over Re this is a straight line in
  $\ln\mathrm{Re}$ (residual 0.03 in $\ln$). This is co-kriging after Kennedy and O'Hagan \cite{kennedy2000} with
  $\rho = 1$ in the logarithmic quantity.
\item \textbf{Finest level:} the series 48/96/144 at Re 3000 is \emph{not} in the asymptotic range. We get
  $d_{96}/d_{144} = 0.58$, while a grid error $\propto n^{-p}$ needs more than $\ln 2/\ln 3 = 0.63$; no positive order
  fits. There are therefore two upper levels. \emph{n144} is computed: $d_{96}(x)$ is scaled by
  $d_{144}/d_{96} \approx 1.7$. \emph{Grid $\to\infty$} is estimated, with Richardson over the last step only and the
  order from the plate test: $d_\infty = d_{144} + (d_{144}-d_{96})/(1.5^{1.33}-1)$. Half the distance between the two
  levels serves as the uncertainty.
\end{enumerate}

\begin{table}[ht]\centering\small
\begin{tabular}{llcccc}\toprule
& & \multicolumn{2}{c}{$\mathrm{Nu}/\mathrm{Nu}_\mathrm{Renon}$} & \multicolumn{2}{c}{$f/f_\mathrm{Renon}$} \\
Type & Re & n144 & grid $\to\infty$ & n144 & grid $\to\infty$ \\\midrule
Gyroid & 1000 & 1.20 & $1.03\pm11\,\%$ & $0.94\pm15\,\%$ & (uncertain) \\
Gyroid & 3000 & 1.29 & $1.03\pm14\,\%$ & $1.06\pm15\,\%$ & (uncertain) \\
Gyroid & 8000 & 1.49 & $1.14\pm16\,\%$ & $1.35\pm15\,\%$ & (uncertain) \\
Diamond & 1000 & 1.14 & $0.97\pm12\,\%$ & $0.98\pm6\,\%$ & (uncertain) \\
Diamond & 3000 & 1.23 & $1.00\pm14\,\%$ & $1.01\pm6\,\%$ & (uncertain) \\
Diamond & 8000 & 1.40 & $1.08\pm16\,\%$ & $1.06\pm6\,\%$ & (uncertain) \\\bottomrule
\end{tabular}
\caption{Surrogate at Pr 1.5 and standard cell. The extrapolated friction cannot be used: it would drop below Renon,
because $f$ does not converge monotonically between the grids.}
\end{table}

\bild{tauscher_surrogat_en.png}{1.0}{Surrogate for gyroid (blue) and diamond (orange). Left, Nu: dotted the coarse level,
dashed n144, solid with band grid $\to\infty$, grey dashed Renon. Centre, $f$ (n144 with band). Right, the
cross-check of the coarse level.}
\lesehilfe{A \emph{surrogate} is a fast stand-in model. It learns from 40 expensive flow
calculations and then gives an instant estimate, with its uncertainty, for any input. Left: the Nusselt number (heat
transfer) versus the \emph{Reynolds number}, a measure of how fast and turbulent the flow is; both axes are
logarithmic. Dots are calculations (circles coarse, squares fine, diamonds finest), lines are the surrogate, the
light band its uncertainty, and the grey dashed line the literature. Centre: the same for the friction factor $f$,
which sets the pressure drop. Right: the cross-check. Each calculation was left out once and predicted from the
others. Dots on the diagonal mean a correct prediction; the deviation is only 1 to 3\,\%.}


\paragraph{What the surrogate says.}
\begin{itemize}
\item Extrapolated to an infinitely fine grid, Nu lands on the Renon correlation (0.97--1.14). The excess of 1.5 at 96
  cells was therefore grid error, not physics. The estimate of 1.2 from only two levels was too high.
\item At equal $d_h$, Nu and $f$ hardly depend on the cell size: for the gyroid, the kernel length scale for $a$ sits at
  its upper bound. A smaller cell gives a smaller $d_h$ but hardly changes the dimensionless numbers.
\item Gyroid and diamond are level for heat transfer. For pressure drop, the gyroid is higher at high Re.
\end{itemize}

\paragraph{Application: does grading pay off?} The preheater of Section~\ref{sec:preheater} was set up again as a
segment model (\texttt{stellarator/vorwaermer\_stufung.py}):
\begin{itemize}
\item Counterflow, two-fluid TPMS; both channel networks share the cell, so $d_h$ is the same on both sides.
\item 8 segments of equal heat duty, each with local temperatures, properties, Pr and LMTD.
\item Nu and $f$ from the surrogate (level n144; Nu additionally at grid $\to\infty$ as a band).
\item \emph{Uniform:} one cell for all segments. \emph{Graded:} each segment chooses its own cell so that
  $P_i + \lambda V_i$ is minimal, with $\lambda$ sampled.
\item Also free: approach, condensate velocity, and number of passes on the warm side.
\end{itemize}

\bild{vorwaermer_stufung_en.png}{0.85}{Net extra power over exchanger volume. Graded (dashed) and uniform coincide. For
comparison, Renon and Renon with Nu $\times$ 1.28 (grading after Oh et al.) in the same model.}
\lesehilfe{Same layout as the net-gain figure above, now with the surrogate. The dashed lines are
exchangers whose cell size changes in eight steps along the flow path (\emph{graded}). The solid lines are
exchangers with the same cells everywhere. They lie practically on top of each other. The band spans from the
finely computed to the extrapolated heat transfer. Grading brings no benefit; the smallest cell is best
everywhere.}


\paragraph{Result.}
\begin{itemize}
\item \textbf{Grading brings nothing.} In every segment the smallest allowed cell wins ($a = 0.75\,a_\mathrm{std}$,
  $d_h = 1.33$\,mm). The graded design is therefore identical to the uniform one. The reason is that pumping power is
  small here (below 1\,MW). Power is lost mainly through the approach at the warm end, and the cheapest way to reduce
  that is more area per volume, i.e.\ fine cells everywhere.
\item \textbf{Volume:} the preheater reaches 95\,\% of the 73.3\,MW upper limit with 1.3\,m$^3$ (n144) to 1.5\,m$^3$
  (grid $\to\infty$). At 1\,m$^3$, CFD and Renon give 66.0--67.6\,MW, a small difference.
\item The literature factor for grading (Nu $\times$ 1.28 \cite{oh2025}) here matches almost exactly the gap between
  Renon and the n144 surrogate. The improvement reported by Oh et al.\ cannot be reproduced at the cell level. It must
  come from effects this model does not contain, such as the distribution in the headers.
\item \textbf{Limits:} the best cell lies at the edge of the surrogate range; finer cells have not been computed. Whether
  the 0.5\,mm wall carries the 3\,MPa pressure difference has not been checked. Manufacturing such fine cells is an open
  question. Grading may pay off where pumping power matters, for example on the steam side of the condenser.
\end{itemize}
